\scanpage{52}
\deleted{\unsure{de}} \uline{canonique de} $\Delta\inserted{\subset F}$ \uline{n'est autre que la translation par} $\underline b$ (cas particulier de (94) en car. 2). \deleted{Donc, pour $\underline b$ fixé, on peut l'exprimer}

Dire que $\sigma$ dans $B$ est une réflexion (\deleted{pas} seult une pseudoréflexion) signifie ici que $\underline b$ est non nul en chaque fibre\deleted{\unsure{s}} (§ 7), ou encore (\unsure{p.c.q.|p.ex.|par ex.} en vertu de (98)) que $D$ est \deleted{\unsure{inversible|partout}} \inserted{section} inversible, ou enfin que $\Delta$ est étale sur $S$.

Dans ce cas, pour $\underline b$ fixé, les $\Delta$ correspondants (i.e. les $q$ correspondantes \ldots) correspondent biunivoquement aux sections du schéma
\[F/\underline b\]

Or \deleted{si on considère l'hom} on aura une suite exacte de schémas en groupes sur $S$
\diagram{\textup{(101)}\quad$\begin{tikzcd}[ampersand replacement=\&,column sep=small]
0\arrow[r]\&(\mathbb Z/2\mathbb Z)_S\arrow[r,swap,"{\varepsilon\mapsto\varepsilon\underline b}"]\&\check{\mathbb V}(L)\arrow[r,"\beta","{x\mapsto x^2+\underline b x}"']\&\check{\mathbb V}(L^{\otimes2})\arrow[r]\&0
\end{tikzcd}$}
% slanted
\margin{Suite exacte d'Artin-Schreyer (en car. 2)}
de sorte que
\[F/\underline b\simeq F\mathbin{\wedge}_{(\check{\mathbb V}(L),\beta)}\check{\mathbb V}(L^{\otimes2})\overset{\mathrm{déf}}=C\tag*{(102)}\]
s'identifie à un torseur \deleted{\unsure{mult}} sous $\check{\mathbb V}(L^{\otimes2})$, et la donnée de $q$ (ou $\Delta$) compatible avec $\underline b$ s'identifie à la donnée d'une section \inserted{$\sigma$} de ce torseur, dont $F$ est un rev. \uline{étale} quadratique $\Delta$ se récupère alors comme l'image inverse par $\sigma$ de ce rev. étale quadratique.
% archive sheet number: 51
